
\documentclass{article}
\usepackage{colortbl}
\usepackage{makecell}
\usepackage{multirow}
\usepackage{supertabular}

\begin{document}

\newcounter{utterance}

\centering \large Interaction Transcript for game `codenames', experiment `abstract', episode 5 with reflexion{-}llama{-}3.1{-}8b{-}instruct.
\vspace{24pt}

{ \footnotesize  \setcounter{utterance}{1}
\setlength{\tabcolsep}{0pt}
\begin{supertabular}{c@{$\;$}|p{.15\linewidth}@{}p{.15\linewidth}p{.15\linewidth}p{.15\linewidth}p{.15\linewidth}p{.15\linewidth}}
    \# & $\;$A & \multicolumn{4}{c}{Game Master} & $\;\:$B\\
    \hline

    \theutterance \stepcounter{utterance}  
    & & & \multicolumn{2}{p{0.3\linewidth}}{
        \cellcolor[rgb]{0.9,0.9,0.9}{
            \makecell[{{p{\linewidth}}}]{
                \texttt{\tiny{[GM$|$GM]}}
                {'hidden': {'team': ['priority'], 'innocent': ['urge', 'reaffirmation', 'hilarity'], 'opponent': ['unfriendliness'], 'assassin': []}, 'revealed': {'team': {'team': [], 'innocent': [], 'opponent': [], 'assassin': []}, 'opponent': {'team': [], 'innocent': [], 'opponent': [], 'assassin': []}}}
            }
        }
    }
    & & \\ \\

    \theutterance \stepcounter{utterance}  
    & & \multicolumn{4}{p{0.6\linewidth}}{
        \cellcolor[rgb]{0.9,0.9,0.9}{
            \makecell[{{p{\linewidth}}}]{
                \texttt{\tiny{[P1$\langle$GM]}}
                \texttt{You have attempted this same task before and did not succeed. These are the plans you wrote after those attempts.} \\
\texttt{Trial \#1: New plan: } \\
\\ 
\texttt{In my previous attempts, I was able to provide valid clues that helped my teammate guess some of our team's words. However, I made a mistake in my last attempt by providing a clue that was too specific and only related to two words, "whim" and "reluctance". This made it difficult for my teammate to guess the remaining words.} \\
\\ 
\texttt{Next time, I will make sure to provide a clue that is more general and relates to multiple words. I will also pay closer attention to the remaining words on the list and try to find a connection that is common to at least three words. For example, I could have provided a clue like "CLUE: abstract" that relates to multiple words, such as "whim", "unbelief", and "ethos", which are all abstract concepts.} \\
\\ 
\texttt{Trial \#2: New plan: } \\
\\ 
\texttt{In my previous attempts, I was able to provide valid clues that helped my teammate guess some of our team's words. However, I made a mistake in my last attempt by providing a clue that was too specific and only related to two words, "whim" and "reluctance". This made it difficult for my teammate to guess the remaining words.} \\
\\ 
\texttt{Next time, I will make sure to provide a clue that is more general and relates to multiple words. I will also pay closer attention to the remaining words on the list and try to find a connection that is common to at least three words. For example, I could have provided a clue like "CLUE: abstract" that relates to multiple words, such as "whim", "unbelief", and "ethos", which are all abstract concepts.} \\
\\ 
\texttt{To achieve this, I will start by looking for common themes or categories among the remaining team words, such as emotions, abstract concepts, or actions. I will then try to find a word that fits into one of these categories and relates to at least three of the remaining team words. This will increase the chances of my teammate guessing the remaining words correctly.} \\
\\ 
\texttt{Trial \#3: New plan: } \\
\\ 
\texttt{In my previous attempts, I was able to provide valid clues that helped my teammate guess some of our team's words. However, I made a mistake in my last attempt by providing a clue that was too specific and only related to two words, "priority" and "urge". This made it difficult for my teammate to guess the remaining words.} \\
\\ 
\texttt{Next time, I will make sure to provide a clue that is more general and relates to multiple words. I will also pay closer attention to the remaining words on the list and try to find a connection that is common to at least three words. For example, I could have provided a clue like "CLUE: concept" that relates to multiple words, such as "priority", "urge", and "unfriendliness", which are all abstract concepts.} \\
\\ 
\texttt{To achieve this, I will start by looking for common themes or categories among the remaining team words, such as emotions, abstract concepts, or actions. I will then try to find a word that fits into one of these categories and relates to at least three of the remaining team words. This will increase the chances of my teammate guessing the remaining words correctly.} \\
\\ 
\\ 
\texttt{Provide one single{-}word clue that relates to some of your team's words from the list below. You can choose to how many and to which words your clue relates to.} \\
\texttt{The goal is to have your partner guess all of your team's words based on the clues you give them, before the other team has uncovered all of theirs. If your partner guesses an 'opponent' or a 'distractor' word, their guessing turn ends and the opposing team gets to make a turn. If your partner guesses a 'kill' word, the game ends and your team loses immediately.} \\
\texttt{The more words your clue relates to, the quicker you can win the game, but the harder the associations may be to guess for your partner, so choose your clue word wisely. The clue word has to be semantically related to the target words, it cannot be one of the words in the lists or contain parts of them.} \\
\\ 
\texttt{Always give your single{-}word clue and your comma{-}separated list of related target words in the following format and make your answers as short as possible, never include any other text than is required in this form:} \\
\\ 
\texttt{CLUE: <WORD>} \\
\texttt{TARGETS: <TARGETS>} \\
\\ 
\texttt{Your team words are: priority.} \\
\texttt{Your opponent's team words are: unfriendliness.} \\
\texttt{Distractor words are: urge, reaffirmation, hilarity.} \\
\texttt{The kill words are: .} \\
            }
        }
    }
    & \\ \\

    \theutterance \stepcounter{utterance}  
    & \multicolumn{4}{p{0.6\linewidth}}{
        \cellcolor[rgb]{0.9,0.9,0.9}{
            \makecell[{{p{\linewidth}}}]{
                \texttt{\tiny{[P1$\rangle$GM]}}
                \texttt{CLUE: abstract} \\
\texttt{TARGETS: priority, unfriendliness} \\
            }
        }
    }
    & & \\ \\

    \theutterance \stepcounter{utterance}  
    & & & \multicolumn{2}{p{0.3\linewidth}}{
        \cellcolor[rgb]{0.9,0.9,0.9}{
            \makecell[{{p{\linewidth}}}]{
                \texttt{\tiny{[GM$|$GM]}}
                \texttt{abstract} \\
            }
        }
    }
    & & \\ \\

    \theutterance \stepcounter{utterance}  
    & & & \multicolumn{2}{p{0.3\linewidth}}{
        \cellcolor[rgb]{0.9,0.9,0.9}{
            \makecell[{{p{\linewidth}}}]{
                \texttt{\tiny{[GM$|$GM]}}
                ['priority', 'unfriendliness']
            }
        }
    }
    & & \\ \\

    \theutterance \stepcounter{utterance}  
    & & & \multicolumn{2}{p{0.3\linewidth}}{
        \cellcolor[rgb]{0.9,0.9,0.9}{
            \makecell[{{p{\linewidth}}}]{
                \texttt{\tiny{[GM$|$GM]}}
                {'word': 'priority', 'assignment': 'team'}
            }
        }
    }
    & & \\ \\

    \theutterance \stepcounter{utterance}  
    & & & \multicolumn{2}{p{0.3\linewidth}}{
        \cellcolor[rgb]{0.9,0.9,0.9}{
            \makecell[{{p{\linewidth}}}]{
                \texttt{\tiny{[GM$|$GM]}}
                {'word': 'unfriendliness', 'assignment': 'opponent'}
            }
        }
    }
    & & \\ \\

    \theutterance \stepcounter{utterance}  
    & & \multicolumn{4}{p{0.6\linewidth}}{
        \cellcolor[rgb]{0.9,0.9,0.9}{
            \makecell[{{p{\linewidth}}}]{
                \texttt{\tiny{[GM$\rangle$P2]}}
                \texttt{You have attempted this same task before and did not succeed. These are the plans you wrote after those attempts.} \\
\texttt{Trial \#1: New plan: In my last attempt, I got stuck in an infinite loop of explaining the association between 'whim' and 'hesitation', which is not a valid response. What I should have done is to provide a concise explanation and a list of words associated with 'hesitation'. I should have avoided repeating the same explanation multiple times and instead focused on providing a clear and well{-}formed response. Next time, I will make sure to provide a single, clear explanation and a list of up to 2 words that best relate or are most closely associated with the word 'hesitation'.} \\
\\ 
\texttt{Trial \#2: New plan: In my last attempt, I failed to achieve the target state of the game because I kept guessing words that were either opponent words or innocent words. What I should have done is to pay closer attention to the list of words provided and focus on guessing words that are closely associated with the given word. I should have also avoided repeating the same list of words multiple times and instead tried to come up with new and more relevant associations. Next time, I will make sure to carefully read the list of words and provide a list of up to 2 words that best relate or are most closely associated with the given word, and I will try to come up with new and more relevant associations each time.} \\
\\ 
\texttt{Trial \#3: New plan: In my last attempt, I failed to provide a well{-}formed and legal response because I included a word ("priority") that was not in the list of provided words. What I should have done is to carefully read the list of words and only include words that were actually provided. I should have also made sure to start my response with the exact phrase "GUESS: " followed by a comma{-}separated list of up to 3 words. Next time, I will make sure to carefully read the list of words and provide a list of up to 3 words that best relate or are most closely associated with the given word, and I will make sure to follow the exact format specified by the Game Master.} \\
\\ 
\\ 
\texttt{Provide a comma{-}separated list of up to 2 words from the following list that best relate or are most closely associated with the word 'abstract'. Always start your list of guess(es) with 'GUESS: ' and do not include any other text in your answer.} \\
\\ 
\texttt{hilarity, reaffirmation, priority, urge, unfriendliness} \\
            }
        }
    }
    & \\ \\

    \theutterance \stepcounter{utterance}  
    & & & \multicolumn{4}{p{0.6\linewidth}}{
        \cellcolor[rgb]{0.9,0.9,0.9}{
            \makecell[{{p{\linewidth}}}]{
                \texttt{\tiny{[GM$\langle$P2]}}
                \texttt{GUESS: concept, idea} \\
            }
        }
    }
     \\ \\

    \theutterance \stepcounter{utterance}  
    & & & \multicolumn{2}{p{0.3\linewidth}}{
        \cellcolor[rgb]{0.9,0.9,0.9}{
            \makecell[{{p{\linewidth}}}]{
                \texttt{\tiny{[GM$|$GM]}}
                {'player': 'Guesser', 'type': 'guess word is hallucination', 'utterance': 'concept, idea', 'message': "Guessed word 'concept' was not listed, you can only guess words provided in the lists.", 'guess': 'concept', 'previous_guesses': [], 'remaining_words': ['hilarity', 'reaffirmation', 'priority', 'urge', 'unfriendliness']}
            }
        }
    }
    & & \\ \\

    \theutterance \stepcounter{utterance}  
    & \multicolumn{4}{p{0.6\linewidth}}{
        \cellcolor[rgb]{0.95,0.95,0.95}{
            \makecell[{{p{\linewidth}}}]{
                \texttt{\tiny{[P2$\rightarrow$P1]}}
                \texttt{Post{-}episode Reflexion (not sent to GM):} \\
\texttt{New plan: } \\
\\ 
\texttt{In my previous attempts, I was able to provide valid clues that helped my teammate guess some of our team's words. However, I made a mistake in my last attempt by providing a clue that was too specific and only related to two words, "priority" and "unfriendliness". This made it difficult for my teammate to guess the remaining words.} \\
\\ 
\texttt{Next time, I will make sure to provide a clue that is more general and relates to multiple words. I will also pay closer attention to the remaining words on the list and try to find a connection that is common to at least three words. For example, I could have provided a clue like "CLUE: concept" that relates to multiple words, such as "priority", "urge", and "unfriendliness", which are all abstract concepts.} \\
\\ 
\texttt{To achieve this, I will start by looking for common themes or categories among the remaining team words, such as emotions, abstract concepts, or actions. I will then try to find a word that fits into one of these categories and relates to at least three of the remaining team words. This will increase the chances of my teammate guessing the remaining words correctly.} \\
            }
        }
    }
    & & \\ \\

    \theutterance \stepcounter{utterance}  
    & & & \multicolumn{4}{p{0.6\linewidth}}{
        \cellcolor[rgb]{0.95,0.95,0.95}{
            \makecell[{{p{\linewidth}}}]{
                \texttt{\tiny{[P1$\rightarrow$P2]}}
                \texttt{Post{-}episode Reflexion (not sent to GM):} \\
\texttt{New plan: In my last attempt, I failed to provide a well{-}formed and legal response because I included two words ("concept" and "idea") that were not a comma{-}separated list of up to 2 words from the provided list. What I should have done is to carefully read the list of words and only include words that were actually provided. I should have also made sure to start my response with the exact phrase "GUESS: " followed by a comma{-}separated list of up to 2 words. Next time, I will make sure to carefully read the list of words and provide a list of up to 2 words that best relate or are most closely associated with the given word, and I will make sure to follow the exact format specified by the Game Master.} \\
            }
        }
    }
     \\ \\

\end{supertabular}
}

\end{document}
